\chapter{Systemy rankingowe jako modele bazowe}

\section{Rola systemów rankingowych w predykcji}

Systemy rankingowe stanowią pierwszy poziom modelowania siły sportowej. Ich zadaniem jest
przekształcenie historii wyników w liczbową ocenę aktualnej siły drużyny lub zawodnika. W kontekście
predykcji przedmeczowej ranking jest użyteczny, ponieważ może zostać obliczony wyłącznie na
podstawie informacji dostępnych przed danym meczem. Po zakończeniu meczu ranking jest aktualizowany
i może wpływać dopiero na kolejne predykcje.

W pracy systemy rankingowe pełnią dwie funkcje. Po pierwsze, są samodzielnymi modelami bazowymi,
które można porównać z benchmarkiem bukmacherskim. Po drugie, generują cechy wejściowe dla
metamodelu uczenia maszynowego. Metamodel nie zastępuje więc rankingów, lecz wykorzystuje je jako
syntetyczne reprezentacje siły stron meczu.

\section{Ranking drużynowy i ranking zawodniczy}

W klasycznym ujęciu ranking może być przypisany drużynie. Takie podejście jest naturalne w wielu
sportach zespołowych, ponieważ nazwa klubu często jest stabilnym nośnikiem informacji o sile. W
profesjonalnym League of Legends założenie to jest jednak ograniczone. Organizacje esportowe
regularnie zmieniają składy, a pojedynczy zawodnicy przenoszą znaczną część informacji o jakości gry
pomiędzy drużynami.

Z tego powodu w pracy porównano dwa poziomy modelowania:

\begin{itemize}
    \item \textbf{ranking drużynowy}, w którym aktualizowana jest ocena organizacji lub drużyny jako całości,
    \item \textbf{ranking zawodniczy}, w którym aktualizowane są oceny graczy występujących w danym meczu, a siła zespołu jest wyznaczana na podstawie składu.
\end{itemize}

Hipoteza badawcza H1 zakłada, że podejście zawodnicze powinno być skuteczniejsze od drużynowego.
Wynika to z faktu, że w LoL skład osobowy jest bardziej bezpośrednim opisem rzeczywistej siły
zespołu niż sama nazwa organizacji.

\section{Zastosowane rodziny rankingów}

\subsection{Elo}

System Elo jest jednym z najprostszych i najbardziej znanych systemów oceny siły graczy
\parencite{elo1978}. Model zakłada, że oczekiwany wynik zależy od różnicy rankingów dwóch stron. Po
meczu ranking zwycięzcy rośnie, a ranking przegranego maleje. Wielkość zmiany zależy od tego, jak
bardzo wynik różnił się od oczekiwania.

Zaletą Elo jest prostota i łatwa interpretacja. Ograniczeniem jest natomiast brak jawnego
modelowania niepewności. Zawodnik z dużą liczbą rozegranych meczów i zawodnik z bardzo krótką
historią mogą mieć podobny rating, mimo że pewność oszacowania ich siły jest różna.

\subsection{Glicko-2}

Glicko-2 rozwija ideę rankingów Elo przez wprowadzenie dodatkowych parametrów opisujących niepewność
oceny \parencite{glickman2012}. Oprócz samego ratingu model uwzględnia rating deviation, czyli miarę
niepewności oszacowania siły. Jest to szczególnie istotne w środowisku esportowym, gdzie część
zawodników ma długą historię występów, a część pojawia się w danych tylko okresowo.

W praktyce oznacza to, że wynik przeciwko dobrze znanemu zawodnikowi może być interpretowany inaczej
niż wynik przeciwko zawodnikowi o bardzo niepewnym poziomie. Ta własność jest korzystna w danych
obejmujących wiele lat, regionów i poziomów rozgrywek.

\subsection{TrueSkill i OpenSkill}

TrueSkill modeluje umiejętność gracza jako rozkład prawdopodobieństwa, a nie pojedynczą wartość
punktową \parencite{trueskill2007}. Podejście to zostało zaprojektowane z myślą o grach
wieloosobowych i zespołowych. OpenSkill jest otwartą implementacją podobnej klasy systemów
rankingowych. W pracy modele te są traktowane jako alternatywne sposoby reprezentacji siły
zawodników i drużyn.

\subsection{Modele porównawcze}

Uzupełniająco wykorzystano także modele z rodziny Plackett--Luce oraz Thurstone--Mosteller. Ich rola
w pracy jest głównie porównawcza: pozwalają sprawdzić, czy przewaga konkretnego systemu rankingowego
wynika z ogólnej idei rankingu, czy ze specyfiki wybranej metody aktualizacji.

\section{Procedura obliczania predykcji rankingowych}

Predykcje rankingowe były generowane chronologicznie. Dla każdego meczu najpierw odczytywano
aktualne rankingi drużyn lub zawodników, następnie obliczano prawdopodobieństwo zwycięstwa pierwszej
strony, a dopiero po tym aktualizowano rankingi na podstawie wyniku. Taka procedura jest zgodna z
zasadą unikania przecieku informacji opisaną w poprzednim rozdziale.

W przypadku rankingów zawodniczych siła strony była wyznaczana na podstawie zawodników występujących
w danym meczu. Dzięki temu transfer zawodnika wpływał na ocenę nowego zespołu, zamiast pozostawać
wyłącznie w historii poprzedniej organizacji.

\section{Dobór parametrów systemów rankingowych}

Same definicje rodzin rankingowych nie przesądzają jeszcze o ich jakości predykcyjnej. Każdy system
wymaga przyjęcia parametrów opisujących tempo aktualizacji, skalę różnic siły lub sposób traktowania
niepewności. Z tego powodu przed wyborem rankingów do dalszych etapów wykonano oddzielny eksperyment
doboru parametrów. Nie był on traktowany jako strojenie końcowego modelu uczenia maszynowego, lecz
jako sprawdzenie, czy słabsze wyniki niektórych rodzin rankingowych nie wynikają wyłącznie z
niekorzystnego ustawienia parametrów domyślnych.

Dobór przeprowadzono w sposób chronologiczny. Okres 2021--2023 wykorzystano jako część walidacyjną
do wyboru konfiguracji, natomiast okres 2024+ jako okres testowy. Dla Elo sprawdzano przede
wszystkim wartość współczynnika aktualizacji $K$, dla Glicko-2 początkowe odchylenie rankingu oraz
tempo wzrostu niepewności, dla TrueSkill parametry $\beta$ i $\tau$, a dla rodzin Plackett--Luce i
Thurstone--Mosteller parametry skali i dryfu. Dla OpenSkill testowano przede wszystkim początkową
wartość niepewności $\sigma$.

Taka procedura pozwala uczciwie porównać rodziny rankingowe po ich podstawowym dostrojeniu.
Jednocześnie zachowuje warunek przedmeczowy: w każdej predykcji ranking był liczony wyłącznie na
podstawie wcześniejszych meczów.

Po wyborze najlepszych konfiguracji predykcje rankingowe zostały wygenerowane ponownie i to one
stanowiły wejście do dalszych eksperymentów metamodelu. Dzięki temu wyniki w kolejnych rozdziałach
są spójne z parametrami raportowanymi w tym rozdziale.

\section{Wyniki po doborze parametrów}

Najważniejsze porównanie rankingów w pracy dotyczy meczów, które można później zestawić również z
rynkiem bukmacherskim. Dzięki temu etap rankingowy pozostaje spójny z końcową próbą ewaluacyjną.
Wyniki w tabeli~\ref{tab:ranking_common_sample_results} pokazują najlepsze konfiguracje wybranych
rodzin rankingowych na okresie testowym 2024+.

W tym rozdziale Player Glicko-2 pełni rolę silnego, prostego baseline'u sportowego. Jest to
odpowiednik prostego modelu odniesienia: nie wykorzystuje pełnego kontekstu historycznego ani
dodatkowego metamodelu, ale pozwala sprawdzić, jak daleko można dojść za pomocą samej
chronologicznej oceny siły zawodników.

\begin{table}[H]
    \centering
    \small
    \caption{Najlepsze konfiguracje rodzin rankingowych na okresie testowym 2024+}
    \label{tab:ranking_common_sample_results}
    \begin{tabular}{p{3.1cm}p{4.1cm}rrrr}
        \toprule
        Rodzina & Najlepsze parametry & AUC & LogLoss & Brier & ECE \\
        \midrule
        Player Glicko-2 & $RD=350$, $vol=0{,}06$, okres $=7$ & 0{,}7511 & 0{,}5941 & 0{,}2040 & 0{,}0443 \\
        Player Elo & $K=48$ & 0{,}7348 & 0{,}6015 & 0{,}2077 & 0{,}0283 \\
        Player Thurstone--Mosteller & $\beta=18{,}75$, $\tau=0{,}05$ & 0{,}7336 & 0{,}6052 & 0{,}2084 & 0{,}0326 \\
        Player Plackett--Luce & $\beta=18{,}75$, $\tau=0{,}05$ & 0{,}7352 & 0{,}6134 & 0{,}2105 & 0{,}0548 \\
        Player TrueSkill & $\beta=4{,}16$, $\tau=0{,}25$ & 0{,}7255 & 0{,}6224 & 0{,}2162 & 0{,}0622 \\
        Player OpenSkill & $\sigma=3{,}50$ & 0{,}7337 & 0{,}6453 & 0{,}2163 & 0{,}0834 \\
        \bottomrule
    \end{tabular}
\end{table}

Wyniki w tabeli~\ref{tab:ranking_common_sample_results} wskazują, że najlepszym pojedynczym systemem
rankingowym jest Player Glicko-2. Model ten osiąga najwyższe AUC oraz najniższy LogLoss spośród
analizowanych rankingów również po doborze parametrów. Oznacza to, że przewaga tej rodziny nie
wynikała wyłącznie z przypadkowego ustawienia parametrów domyślnych.

\begin{figure}[H]
    \centering
    \includegraphics[width=0.9\textwidth]{baseline_point5/baseline_player_vs_team_auc_2020.png}
    \caption{Porównanie jakości rankingów zawodniczych i drużynowych według AUC.}
    \label{fig:ranking_player_vs_team_auc}
\end{figure}

\Cref{fig:ranking_player_vs_team_auc} pokazuje, że ranking oparty na zawodnikach jest zasadniczo
silniejszy niż ranking oparty wyłącznie na drużynie. Jest to zgodne ze specyfiką League of Legends:
historia organizacji nie zawsze dobrze opisuje aktualny skład, natomiast informacje o zawodnikach
lepiej przenoszą się pomiędzy sezonami i transferami.

\section{Porównanie ratingów z benchmarkiem rynkowym}

W porównaniu z benchmarkiem rynkowym Player Glicko-2 osiąga lepszy AUC i LogLoss niż kurs otwarcia
oraz wynik zbliżony do kursu zamknięcia. Nie oznacza to jednak, że surowy ranking jest kompletnym
modelem predykcyjnym. Rynek bukmacherski pozostaje silnie skalibrowanym benchmarkiem, a surowy
ranking ma gorsze ECE niż kursy rynkowe.

\begin{figure}[H]
    \centering
    \includegraphics[width=0.9\textwidth]{baseline_point5/baseline_auc_ranking_common_2020.png}
    \caption{Porównanie AUC rankingów i benchmarków rynkowych na wspólnej próbie.}
    \label{fig:ranking_auc_common}
\end{figure}

\begin{figure}[H]
    \centering
    \includegraphics[width=0.9\textwidth]{baseline_point5/baseline_calibration_common_2020.png}
    \caption{Kalibracja wybranych rankingów i benchmarków rynkowych na wspólnej próbie.}
    \label{fig:ranking_calibration_common}
\end{figure}

Porównanie AUC i wykresów kalibracji prowadzi do dwóch wniosków. Po pierwsze, ranking zawodniczy
zawiera silny sygnał sportowy, ponieważ dobrze porządkuje mecze według prawdopodobieństwa
zwycięstwa. Po drugie, sama wartość ratingowa nie wystarcza do uzyskania najlepszej jakości
probabilistycznej, ponieważ kalibracja prawdopodobieństw może być gorsza niż w przypadku rynku.
Zamiast pełnych krzywych ROC pokazano zestawienie AUC, ponieważ krzywe najlepszych wariantów były
bardzo zbliżone i na jednym wykresie praktycznie się nakładały.

\section{Znaczenie wyników rankingowych dla dalszego modelowania}

Wyniki systemów rankingowych uzasadniają budowę metamodelu. Skoro Player Glicko-2 jest silnym
pojedynczym sygnałem, może stanowić podstawę modelu wyższego poziomu. Jednocześnie jego ograniczenia
kalibracyjne wskazują, że warto połączyć różne informacje rankingowe, niepewność oraz historyczny
kontekst formy w jednym modelu probabilistycznym.

Najważniejszy wniosek z tego rozdziału jest zatem dwojaki. Po pierwsze, ranking zawodniczy jest
lepszym sposobem reprezentacji siły w League of Legends niż ranking drużynowy. Po drugie, najlepszy
pojedynczy ranking nie kończy procesu modelowania, lecz stanowi punkt wyjścia dla metamodelu
opisanego w kolejnym rozdziale.
